% =====================================================================
% Appendix D — Reproducibility (hyperparameters + per-epoch trajectory)
% =====================================================================

\section{Reproducibility}
\label{app:repro}

\paragraph{Hyperparameters.}

\begin{table}[h]
\centering\small
\caption{Hyperparameters for the baseline (\S\ref{sec:baseline}) and the encoder-decoder variants (\S\ref{sec:archsearch}).}
\label{tab:hparams}
\begin{tabular}{@{}lll@{}}
\toprule
& Baseline & Encoder-decoder (both variants) \\
\midrule
History length $S$            & 24                       & 24 \\
Future length $H$              & 24                       & 24 \\
Spatial token grid             & $8\times 8$ (P=64)       & $8\times 8$ (P=64) \\
Embedding dim $D$              & 128                      & 128 \\
Encoder layers                 & 4                        & 4 \\
Decoder layers                 & ---                      & 2 \\
Attention heads                & 4                        & 4 \\
Dropout                        & 0.1                      & 0.1 \\
Optimizer                      & AdamW \citep{loshchilov2019adamw} & AdamW \\
Initial LR                     & $3\times 10^{-4}$        & $3\times 10^{-4}$ \\
Weight decay                   & $10^{-2}$                & $10^{-2}$ \\
Schedule                       & cosine warm-restarts \citep{loshchilov2017sgdr} & cosine warm-restarts \\
Batch size                     & 16                       & 16 \\
Epochs (planned / actual)      & 14 / 14                  & 24 / 12 (v1), 24 / 13 (v2) \\
Loss                           & MSE in z-score space     & MSE in z-score space \\
Mixed precision                & PyTorch AMP fp16          & PyTorch AMP fp16 \\
Random seeds                   & not set                  & not set \\
\bottomrule
\end{tabular}
\end{table}

\paragraph{Per-epoch validation MAPE trajectory.}

\begin{table}[h]
\centering\small
\caption{Validation MAPE per training epoch (lower is better). Encoder-decoder variants used a chained \texttt{--resume} workflow that reset the LR scheduler at every chain boundary; the kink at epoch 7 is visible in Figure~\ref{fig:v1v2}.}
\label{tab:trajectory}
\begin{tabular}{@{}rrrr@{}}
\toprule
Epoch & Baseline (val 2021) & v1 (val 2022) & v2 (val 2022, undertrained) \\
\midrule
0  & 11.22\,\% & 10.08\,\% & --- \\
4  & 9.16\,\%  & 8.70\,\%  & --- \\
6  & 8.76\,\%  & \textbf{8.63\,\%} ($\star$) & --- \\
8  & ---       & 9.95\,\%  & 8.95\,\% (LR reset to 1e-3) \\
13 & \textbf{6.92\,\%} & 10.50\,\% & \textbf{8.72\,\%} ($\star$) \\
\bottomrule
\multicolumn{4}{l}{\footnotesize ($\star$) = best-of-run; both encoder-decoder variants stalled after the LR-scheduler reset.} \\
\end{tabular}
\end{table}

\paragraph{Random seeds caveat.} The training pipeline does not set \texttt{torch.manual\_seed} / \texttt{np.random.seed} / \texttt{torch.use\_deterministic\_algorithms}. The reported test MAPE numbers (5.24\,\% baseline, 6.82\,\% encoder-decoder, 4.33\,\% / 4.21\,\% ensembles) are pinned to the specific checkpoints in \texttt{runs/cnn\_transformer\_baseline/checkpoints/best.pt} and \texttt{runs/cnn\_encoder\_decoder/checkpoints/best.pt}, both bundled with the public release; they are not bit-reproducible across re-training runs.

\paragraph{Validation script.} The deployment-validation experiment in \S\ref{sec:deploy:validation} is executable as a single Python script:
\begin{lstlisting}[language=Python, basicstyle=\ttfamily\footnotesize]
$ python scripts/validation/reproduce_dec30_2022.py
# fetches HRRR for 2022-12-29 (history) + 2022-12-30 (future analyses)
# runs deployed pipeline, compares to the cluster's stored prediction
# expected output:
#   live MAPE       : 6.41 %
#   cluster MAPE    : 6.54 %
#   max abs diff    : 94 MW (3.6 %)
\end{lstlisting}

\paragraph{Compute.}
Training was done on NVIDIA A100-80\,GB / A100-40\,GB / P100 GPUs. Each baseline epoch takes $\sim$45 minutes; the full 14-epoch training is $\sim$10 hours of GPU time. Encoder-decoder epochs are slightly slower ($\sim$60 min/epoch) due to the additional decoder stack. Inference is CPU-feasible: HuggingFace Spaces' free CPU tier (2 vCPU, 16\,GB RAM) serves a forward pass in $\sim$5\,s when warm.

\paragraph{Build instructions.}
The arXiv source is at \href{https://github.com/jeffliulab/real-time-power-predict/tree/main/report/arxiv}{\texttt{report/arxiv/}} in the public code repository. To rebuild:
\begin{lstlisting}[basicstyle=\ttfamily\footnotesize]
$ cd report/arxiv
$ latexmk -xelatex -interaction=nonstopmode -halt-on-error paper.tex
\end{lstlisting}

The figures shipped with the paper are deterministically regenerable via the scripts in \texttt{scripts/figures/}; the live-deployment numbers in Table~\ref{tab:deployment_overall} and Figures~\ref{fig:deploy_per_zone}--\ref{fig:training_vs_deploy} pull from the data-repo URL and will reflect the most recent cron build at compile time. The version of the paper distributed with this repository pins to the 2026-05-05 cron snapshot; subsequent builds will pick up newer windows.
